\section{Verification}
\label{sec:results}

The design is verified at two levels that meet in the middle at the byte-level
protocol. On the hardware side, VHDL testbenches driven through the UART
interface exercise the processing element in isolation, the full protocol
(including checksum rejection and register read-back), and the systolic
matrix--vector mapping end to end. On the host side, the software model is used
as a golden reference: it re-implements the PE arithmetic and the protocol FSM,
so a discrepancy between the two implementations surfaces immediately.

Because the model is byte-accurate, every accelerated kernel is checked against
an independent host computation without any hardware: element-wise arithmetic,
the dot-product reduction, the tiled systolic matrix--vector product, prefix
scans and convolution. The robustness features are exercised too---a fault-
injecting variant of the model forces the retransmission path, confirming that a
rejected transfer is recovered transparently. The same known-answer suite is
reused as a bring-up check on real hardware once the board is programmed.

Table~\ref{tab:verif} summarises the verification surface. The bitstream flow
itself is scripted for batch synthesis but is out of the scope of this logical
report.

\begin{table}[t]
  \centering
  \caption{Verification surface.}
  \label{tab:verif}
  \small
  \begin{tabular}{@{}lp{5.4cm}@{}}
    \toprule
    Level & Coverage \\
    \midrule
    RTL (simulation) & PE ALU; full protocol with checksum; systolic
                       matrix--vector \\
    Model (host)     & Protocol framing; PE arithmetic; all kernels vs.
                       reference; retransmission \\
    Cross-check      & Model is the golden reference for the RTL (shared
                       semantics) \\
    \bottomrule
  \end{tabular}
\end{table}
